\section{Q3. O que seria um projeto Open source? Como empresas podem adotar tais tecnologias e o que isso acarreta?}

Um projeto Open Source é aquele feito colaborativamente por voluntários (ou incomumente, por profissionais pagos) e com o código fonte não compilado disponível para o público geral (por isso o nome código aberto), normalmente por sites como Github. Isso permite que quaisquer pessoas ofereçam sugestões de mudanças, que podem ou não ser acatadas e anexadas ao projeto principal. Projetos Open Source não precisam de pagamento de licença para serem utilizados, podendo ser instalados, distribuídos e modificados livremente.
Empresas podem adotar estas tecnologias simplesmente as baixando e utilizando (juntamente com um treinamento adequado). Ademais, uma vez que projetos open Source podem ser modificados livremente, é plenamente adequado que uma empresa modifique o código para que o programa melhor se adeque à realidade da companhia.